\documentclass[11pt,a4paper]{article}
\usepackage{fontspec,amsmath,amssymb,amsthm,microtype}
\setmainfont{lmroman10-regular.otf}[BoldFont=lmroman10-bold.otf,ItalicFont=lmroman10-italic.otf,BoldItalicFont=lmroman10-bolditalic.otf]
\setmonofont{lmmono10-regular.otf}[Scale=MatchLowercase]
\usepackage[margin=25mm]{geometry}
\usepackage[unicode,hidelinks]{hyperref}
\setlength{\emergencystretch}{3em}
\setlength{\parskip}{0.4em}
\setlength{\parindent}{0pt}
\allowdisplaybreaks[2]
\raggedbottom
\title{\texorpdfstring{The sharp deterministic contraction}{The sharp deterministic contraction}}
\author{LangevinLean mathematical exposition}
\date{}
\begin{document}\maketitle
\noindent Source module: \texttt{Unadjusted\allowbreak{}Langevin/\allowbreak{}Sharp\allowbreak{}Contraction.\allowbreak{}lean}. The accompanying HTML pairs this mathematical exposition with the exact Lean statements and records the scope of the formalization.


The preceding square-root factor uses strong monotonicity and Lipschitz continuity separately. The quadratic lower and upper models contain more information. We combine them into the interpolation inequality, then interpret that inequality as a ball containing the gradient difference. This yields the two sharp factors \(1-mh\) and \(Lh-1\), with the change of regime at \(h=2/(m+L)\).



The proof separates the two regimes for the maximum contraction factor \(\max(|1-mh|,|1-Lh|)\). Each scalar bound follows after the strong monotonicity and cocoercivity estimates have supplied the needed inner-product control. The final norm theorem removes the square by nonnegativity, retaining the precise step-regime hypotheses.


\section{\texorpdfstring{Cocoercivity from two quadratic models}{Cocoercivity from two quadratic models}}


Let \(H\) be a real inner-product space, \(f:H\to\mathbb R\), \(g:H\to H\), and \(K\ge0\). Suppose that for all \(a,b\in H\),

\[f(a)+\langle g(a),b-a\rangle\le f(b)\le f(a)+\langle g(a),b-a\rangle+\frac K2\|b-a\|^2.\]

Then for all \(x,y\),

\[\|g(x)-g(y)\|^2\le K\langle g(x)-g(y),x-y\rangle.\]

The proof includes \(K=0\), and does not separately assume that \(g\) is a gradient.



Fix \(a,b\), \(t\ge0\), and set \(v=g(a)-g(b)\), \(z=a-tv\). Compare the lower model based at \(b\) with the upper model based at \(a\), both evaluated at \(z\):

\[f(b)+\langle g(b),z-b\rangle\le f(z)\le f(a)+\langle g(a),z-a\rangle+\frac K2\|z-a\|^2.\]

The displacements are \(z-a=-tv\) and \(z-b=(a-b)-tv\). Substitute them into the two models. Since \(t\ge0\), the quadratic term is \(\|z-a\|^2=t^2\|v\|^2\).



The difference of the linear terms satisfies

\[\langle g(a),v\rangle-\langle g(b),v\rangle=\|v\|^2.\]

Eliminating \(f(z)\) therefore gives

\[\left(t-\frac K2t^2\right)\|g(a)-g(b)\|^2\le f(a)-f(b)-\langle g(b),a-b\rangle.\]

Apply this estimate to \((a,b)=(x,y)\) and \((y,x)\), with the same \(t\ge0\). Adding will cancel the function values.



The two gradient norms are equal and the reversed displacement changes the sign of its inner product. Thus, writing \(A=\|g(x)-g(y)\|^2\) and \(B=\langle g(x)-g(y),x-y\rangle\), we have for every \(t\ge0\),

\[(2t-Kt^2)A\le B.\]

First consider \(K=0\). Taking \(t=1\) gives \(2A\le B\). As \(A\ge0\), this also shows \(B\ge0\).



If the desired conclusion \(A\le0\) were false, then \(A>0\). We could choose

\[t=\frac{B+1}{2A}\ge0.\]

The estimate \(2tA\le B\) would then force \(B+1\le B\).



To justify that substitution, \(A>0\) implies \(\|g(x)-g(y)\|\ne0\), so its square can be cancelled from the denominator. Hence indeed

\[2\frac{B+1}{2A}A=B+1.\]

The contradiction \(B+1\le B\) excludes \(A>0\). Thus \(A=0\), proving the claim when \(K=0\).



Now suppose \(K>0\). Choose \(t=K^{-1}\) in the common estimate:

\[\frac AK\le B.\]

Multiplying by \(K\) gives \(A\le KB\), as required.



{\raggedright\noindent Lean declarations: \texttt{cocoercive\_\allowbreak{}of\_\allowbreak{}quadratic\_\allowbreak{}models}.\par}

\section{\texorpdfstring{The strong-\allowbreak{}smooth interpolation inequality}{The strong-smooth interpolation inequality}}


Let \(H\) be a real Hilbert space, \(f:H\to\mathbb R\) differentiable and \(m\)-strongly convex, and let \(\nabla f\) be \(L\)-Lipschitz, where \(L\ge0\) and \(m\le L\). Then for every \(x,y\),

\[mL\|x-y\|^2+\|\nabla f(x)-\nabla f(y)\|^2\le(m+L)\langle\nabla f(x)-\nabla f(y),x-y\rangle.\]

The inequality holds also at the boundary \(m=L\).



Subtract the strong-convexity quadratic and its gradient:

\[F(v)=f(v)-\frac m2\|v\|^2,\qquad G(v)=\nabla f(v)-mv.\]

The quadratic lower model for \(f\) becomes an ordinary supporting plane for \(F\):

\[F(a)+\langle G(a),b-a\rangle\le F(b).\]

This follows by expanding the squared distance and collecting the terms involving \(m\).



Apply the smooth upper model to \(f\), and substitute the same definitions of \(F\) and \(G\).



Expanding the quadratic terms shows that the upper curvature decreases from \(L\) to \(L-m\):

\[F(b)\le F(a)+\langle G(a),b-a\rangle+\frac{L-m}{2}\|b-a\|^2.\]

Since \(L-m\ge0\), cocoercivity for these two models gives

\[\|G(x)-G(y)\|^2\le(L-m)\langle G(x)-G(y),x-y\rangle.\]

Write \(d=x-y\) and \(v=\nabla f(x)-\nabla f(y)\). The shifted gradient difference is exactly

\[G(x)-G(y)=v-md.\]

Expand the cocoercivity inequality using this identity:

\[\|v\|^2-2m\langle v,d\rangle+m^2\|d\|^2\le(L-m)\bigl(\langle v,d\rangle-m\langle d,d\rangle\bigr).\]

Use \(\langle d,d\rangle=\|d\|^2\) and collect terms. This gives

\[\|v\|^2+mL\|d\|^2\le(m+L)\langle v,d\rangle.\]

When \(m=L\), the cocoercivity constant above is zero, forcing \(G(x)=G(y)\); consequently \(\nabla f(x)-\nabla f(y)=m(x-y)\). Thus no division by \(L-m\) or limiting argument is needed.



{\raggedright\noindent Lean declarations: \texttt{strong\_\allowbreak{}smooth\_\allowbreak{}gradient\_\allowbreak{}interpolation}.\par}

\section{\texorpdfstring{The gradient difference lies in a ball}{The gradient difference lies in a ball}}


Let \(H\) be a real inner-product space, \(x,y,g_x,g_y\in H\), and \(m,L\in\mathbb R\) with \(m\le L\). Suppose

\[mL\|x-y\|^2+\|g_x-g_y\|^2\le(m+L)\langle g_x-g_y,x-y\rangle.\]

Then

\[\left\|(g_x-g_y)-\frac{m+L}{2}(x-y)\right\|\le\frac{L-m}{2}\|x-y\|.\]

The proposed radius is nonnegative because \(m\le L\). Both sides of the desired norm bound are therefore nonnegative, so it is enough to prove the corresponding squared inequality.



Set \(c=(m+L)/2\), \(r=(L-m)/2\), \(d=x-y\), and \(v=g_x-g_y\). Completing the square in the interpolation inequality gives

\[\begin{aligned}\|v-cd\|^2&=\|v\|^2-(m+L)\langle v,d\rangle+c^2\|d\|^2\\&\le(c^2-mL)\|d\|^2=r^2\|d\|^2.\end{aligned}\]

Taking nonnegative square roots proves the claim. If \(m=L\), the radius is zero and \(v=md\).



{\raggedright\noindent Lean declarations: \texttt{gradient\_\allowbreak{}difference\_\allowbreak{}centered\_\allowbreak{}bound}.\par}

\section{\texorpdfstring{The centered Euler-\allowbreak{}step bound}{The centered Euler-step bound}}


Let \(H\) be a real inner-product space, \(m,L,h\in\mathbb R\) with \(m\le L\) and \(h\ge0\), and let \(x,y,g_x,g_y\in H\) satisfy

\[mL\|x-y\|^2+\|g_x-g_y\|^2\le(m+L)\langle g_x-g_y,x-y\rangle.\]

Then

\[\|(x-hg_x)-(y-hg_y)\|\le\left(\left|1-h\frac{m+L}{2}\right|+h\frac{L-m}{2}\right)\|x-y\|.\]

The preceding ball bound gives \(\|v-cd\|\le r\|d\|\), where \(d=x-y\), \(v=g_x-g_y\), \(c=(m+L)/2\), and \(r=(L-m)/2\).



Decompose the Euler difference into its centered part and its deviation:

\[d-hv=(1-hc)d-h(v-cd).\]

The triangle inequality and homogeneity of the norm give

\[\|d-hv\|\le|1-hc|\,\|d\|+h\|v-cd\|,\]

where \(|-h|=h\) because \(h\ge0\).



Multiply the ball bound by \(h\ge0\) and substitute it:

\[\|d-hv\|\le\bigl(|1-hc|+hr\bigr)\|d\|.\]

For \(m=L\), this reduces to the exact scalar factor \(|1-mh|\).



{\raggedright\noindent Lean declarations: \texttt{gradient\_\allowbreak{}step\_\allowbreak{}centered\_\allowbreak{}contract}.\par}

\section{\texorpdfstring{The sharp small-\allowbreak{}step factor}{The sharp small-step factor}}


Let \(H\) be a real inner-product space, \(x,y,g_x,g_y\in H\), and \(m,L,h\in\mathbb R\). Suppose \(m\le L\), \(m+L>0\), \(0\le h\le2/(m+L)\), and

\[mL\|x-y\|^2+\|g_x-g_y\|^2\le(m+L)\langle g_x-g_y,x-y\rangle.\]

Then

\[\|(x-hg_x)-(y-hg_y)\|\le(1-mh)\|x-y\|.\]

Apply the centered bound. The step-size assumption and \(m+L>0\) imply \(h(m+L)\le2\), hence \(1-h(m+L)/2\ge0\).



Remove the absolute value with this sign and simplify its coefficient:

\[\left|1-h\frac{m+L}{2}\right|+h\frac{L-m}{2}=1-h\frac{m+L}{2}+h\frac{L-m}{2}=1-mh.\]

{\raggedright\noindent Lean declarations: \texttt{gradient\_\allowbreak{}step\_\allowbreak{}sharp\_\allowbreak{}contract\_\allowbreak{}small}.\par}

\section{\texorpdfstring{The sharp large-\allowbreak{}step factor}{The sharp large-step factor}}


Let \(H\) be a real inner-product space, \(x,y,g_x,g_y\in H\), and \(m,L,h\in\mathbb R\). Suppose \(m\le L\), \(m+L>0\), \(h\ge0\), \(2/(m+L)\le h\), and

\[mL\|x-y\|^2+\|g_x-g_y\|^2\le(m+L)\langle g_x-g_y,x-y\rangle.\]

Then

\[\|(x-hg_x)-(y-hg_y)\|\le(Lh-1)\|x-y\|.\]

This is a distance bound; it is a strict contraction only when its coefficient is below one.



Apply the centered estimate. Now the step-size assumption gives \(2\le h(m+L)\), so \(1-h(m+L)/2\le0\).



Replace the absolute value by the negative of its argument. The coefficient becomes

\[\left|1-h\frac{m+L}{2}\right|+h\frac{L-m}{2}=h\frac{m+L}{2}-1+h\frac{L-m}{2}=Lh-1.\]

{\raggedright\noindent Lean declarations: \texttt{gradient\_\allowbreak{}step\_\allowbreak{}sharp\_\allowbreak{}contract\_\allowbreak{}large}.\par}

\section{\texorpdfstring{Small-\allowbreak{}step contraction for a potential}{Small-step contraction for a potential}}


Let \(f:H\to\mathbb R\) be differentiable and \(m\)-strongly convex on a real Hilbert space, with an \(L\)-Lipschitz gradient, \(L\ge0\), \(m\le L\), and \(m+L>0\). If \(0\le h\le2/(m+L)\), then for every \(x,y\),

\[\|(x-h\nabla f(x))-(y-h\nabla f(y))\|\le(1-mh)\|x-y\|.\]

The quadratic models imply the strong-smooth interpolation inequality for \(g_x=\nabla f(x)\) and \(g_y=\nabla f(y)\). Apply the small-step Euler bound to these two vectors. Its hypotheses are exactly the stated curvature and step-size assumptions.



{\raggedright\noindent Lean declarations: \texttt{strong\_\allowbreak{}smooth\_\allowbreak{}gradient\_\allowbreak{}step\_\allowbreak{}contract}.\par}

\section{\texorpdfstring{Large-\allowbreak{}step contraction for a potential}{Large-step contraction for a potential}}


Let \(H\) be a real Hilbert space and \(f:H\to\mathbb R\) differentiable and \(m\)-strongly convex, with an \(L\)-Lipschitz gradient. Suppose \(L\ge0\), \(m\le L\), \(m+L>0\), \(h\ge0\), and \(2/(m+L)\le h\). Then for every \(x,y\in H\),

\[\|(x-h\nabla f(x))-(y-h\nabla f(y))\|\le(Lh-1)\|x-y\|.\]

Use the potential's interpolation inequality in the large-step Euler bound. This gives the second sharp factor \(Lh-1\).



{\raggedright\noindent Lean declarations: \texttt{strong\_\allowbreak{}smooth\_\allowbreak{}gradient\_\allowbreak{}step\_\allowbreak{}contract\_\allowbreak{}large}.\par}

\section{\texorpdfstring{The endpoint-\allowbreak{}factor form}{The endpoint-factor form}}


Let \(f\) be differentiable and \(m\)-strongly convex on a real Hilbert space, with an \(L\)-Lipschitz gradient, \(L\ge0\), \(m\le L\), and \(m+L>0\). For every \(h\ge0\) and every \(x,y\),

\[\|(x-h\nabla f(x))-(y-h\nabla f(y))\|\le\max\{|1-mh|,|1-Lh|\}\,\|x-y\|.\]

If \(0<m\le L\), this factor is below one for \(0<h<2/L\). At \(m=L\), the two endpoint factors coincide.



Split according to whether \(h\le2/(m+L)\). In the first case the small-step theorem bounds the distance by \((1-mh)\|x-y\|\).



Since \(1-mh\le|1-mh|\le\max\{|1-mh|,|1-Lh|\}\), multiplying by \(\|x-y\|\ge0\) gives the desired bound in this case.



In the other case \(h>2/(m+L)\), so the large-step theorem bounds the distance by \((Lh-1)\|x-y\|\).



The elementary absolute-value inequality yields

\[Lh-1=-(1-Lh)\le|1-Lh|.\]

Finally \(|1-Lh|\le\max\{|1-mh|,|1-Lh|\}\). Multiply by the nonnegative distance to finish the second case.



{\raggedright\noindent Lean declarations: \texttt{strong\_\allowbreak{}smooth\_\allowbreak{}gradient\_\allowbreak{}step\_\allowbreak{}contract\_\allowbreak{}max}.\par}

\section{\texorpdfstring{Sharp shared-\allowbreak{}noise contraction: small steps}{Sharp shared-noise contraction: small steps}}


Let \(f:E_d\to\mathbb R\) be differentiable and \(m\)-strongly convex, with an \(L\)-Lipschitz gradient, \(L\ge0\), \(m\le L\), and \(m+L>0\). If \(0\le h\le2/(m+L)\), then for all starting states \(x,y\) and every common noise vector \(z\),

\[\|T_h^z(x)-T_h^z(y)\|\le(1-mh)\|x-y\|.\]

The common noise cancels exactly, leaving the distance between the gradient steps from \(x\) and \(y\).



Apply the potential-based small-step estimate. It gives the factor \(1-mh\) independently of the value or distribution of \(z\).



{\raggedright\noindent Lean declarations: \texttt{ula\_\allowbreak{}shared\_\allowbreak{}noise\_\allowbreak{}contract\_\allowbreak{}sharp}.\par}

\section{\texorpdfstring{Sharp shared-\allowbreak{}noise contraction: large steps}{Sharp shared-noise contraction: large steps}}


Let \(f:E_d\to\mathbb R\) be differentiable and \(m\)-strongly convex, with an \(L\)-Lipschitz gradient, \(L\ge0\), \(m\le L\), and \(m+L>0\). If \(h\ge0\) and \(2/(m+L)\le h\), then for every \(x,y,z\in E_d\),

\[\|T_h^z(x)-T_h^z(y)\|\le(Lh-1)\|x-y\|.\]

Cancel the shared noise, reducing the two ULA updates to their deterministic gradient-step difference.



Apply the potential-based large-step estimate. This yields the factor \(Lh-1\) for every common noise vector.



{\raggedright\noindent Lean declarations: \texttt{ula\_\allowbreak{}shared\_\allowbreak{}noise\_\allowbreak{}contract\_\allowbreak{}sharp\_\allowbreak{}large}.\par}

\end{document}
